\chapter{Noetherian Rings}

\begin{proposition}{7.1}
	Any homomorphic image of Noetherian ring is Noetherian.
\end{proposition}

\begin{proposition}{7.2}
	Let $A$ be a subring of $B$; soppose that $A$ is Noetherian and that $B$ is finitely generated as an $A$-module. Then $B$ is Noetherian.
\end{proposition}

\begin{theorem}{7.5}[\textcolor{red}{Hilbert's basis theorem}]
	If $A$ is Noetherian, then the polynomial ring $A[x]$ is Noetherian.
\end{theorem}

\begin{proposition}{7.8}
	Let $A\subseteq B\subseteq C$ be rings. Suppose that $A$ is Noetherian, that $C$ is finitely generated as an $A$-algebra and that $C$ is either (i) finitely generated as a $B$-module or (ii) integral over $B$. Then $B$ is finitely generated as an $A$-algebra.
\end{proposition}

\begin{proposition}{7.9}[\textcolor{red}{Zariski's lemma}]
	Let $k$ be a field, $E$ a finitely generated $k$-algebra. If $E$ is a field then it is a finite algebraic extension of $k$.
\end{proposition}

\section*{PRIMARY DECOMPOSITION IN NOETHERIAN RINGS}

\begin{proposition}{7.11}
	In a Noetherian ring every ideal is a finite intersection of irreducible ideals.
\end{proposition}

\begin{proposition}{7.12}
	In a Noetherian ring every irreducible ideal is primary.
\end{proposition}

\begin{proposition}{7.14}
	In a Noetherian ring every ideal is contains a power of its radical.
\end{proposition}

\begin{proposition}{7.15}
	In a Noetherian ring the nilradical is nilpotent.
\end{proposition}

\begin{proposition}{7.16}
	Let $A$ be a Noetherian ring, $m$ a maximal ideal of $A$, $\mathfrak{q}$ any ideal of $A$. Then the following are equivalent:
	\begin{enumerate}[label=\textup{(\roman*)}]
		\item $\mathfrak{q}$ is $\mathfrak{m}$-primary;
		\item $\sqrt{\mathfrak{q}}=\mathfrak{m}$;
		\item $\mathfrak{m}^n\subseteq\mathfrak{q}\subseteq\mathfrak{m}$ for some $n>0$.
	\end{enumerate}
\end{proposition}

\begin{proposition}{7.17}
	Let $\mathfrak{a}\neq(1)$ be an ideal in a Noetherian ring. Then the prime ideals which belong to $\mathfrak{a}$ are precisely the prime ideals which occor in the set of ideals $(\mathfrak{a}:x)$.
\end{proposition}